% Compilen desde esta carpeta o desde la raíz del repositorio.
\makeatletter
\def\input@path{{../../preamble/}{book/preamble/}}
\makeatother
\input{preamble.tex}

\usepackage[spanish]{babel}
\usepackage{lmodern}
\usepackage{titlesec}
\geometry{a4paper,margin=1.6cm}
\linespread{1.10}

\definecolor{HeadingColor}{HTML}{243B53}
\definecolor{InstructionColor}{HTML}{704214}
\definecolor{ConceptColor}{HTML}{205C50}
\definecolor{LinkColor}{HTML}{0047AB}
\hypersetup{allcolors=LinkColor}
\newcommand{\guidelink}[2]{\href{#1}{\textbf{\underline{#2}}}}
\newcommand{\instruction}[1]{\textcolor{InstructionColor}{\textbf{#1}}}
\newcommand{\concept}[1]{\textcolor{ConceptColor}{\textbf{#1}}}
\newcommand{\checkbox}{%
    \begingroup
    \setlength{\fboxrule}{0.4pt}%
    \setlength{\fboxsep}{2pt}%
    \fbox{\rule{0pt}{0.8ex}\rule{0.8ex}{0pt}}%
    \endgroup}

\setlength{\parskip}{4pt}
\setlength{\tabcolsep}{5pt}
\renewcommand{\arraystretch}{1.15}
\setlist{leftmargin=0.65cm,topsep=4pt,partopsep=0pt,parsep=0pt,itemsep=3pt}
\titleformat{\section}{\large\bfseries\color{HeadingColor}}{}{0pt}{}
\titlespacing*{\section}{0pt}{10pt}{3pt}
\newcolumntype{Y}{>{\raggedright\arraybackslash}X}

\hypersetup{pdftitle={C5--C6. Continúen la construcción y documenten el avance},pdfsubject={Emprendimiento e Innovación, Grado 10, Periodo 2}}

\begin{document}

\begin{center}
{\Large\bfseries\color{HeadingColor} C5--C6. Continúen la construcción y documenten el avance\par}
\vspace{3pt}
Emprendimiento e Innovación \enspace|\enspace Grado 10 \enspace|\enspace Periodo 2\\
Guía 3 de 5
\end{center}

\section*{Objetivo y punto de partida}
\concept{C5.} Aplica principios básicos de diseño, funcionalidad, accesibilidad o experiencia de usuario según el proyecto.\\
\concept{C6.} Organiza recursos, roles y tiempos para avanzar en la construcción de la propuesta.

Recuperen R1, V1 y PL1 (P2, C3--C4), D1 y los H pertinentes (P1). Continúen el mismo proceso: \textbf{C4 materializa la primera versión; C5 aplica y justifica decisiones en el prototipo; C6 organiza y ejecuta su construcción}. Desarrollen C5 y C6 de manera articulada: una decisión puede cambiar las tareas y el avance puede exigir revisar una decisión.

\section*{Documento acumulativo e integridad de la evidencia}
Continúen el \textbf{documento acumulativo en LaTeX} en \guidelink{https://www.overleaf.com/}{Overleaf} o su editor habitual; entreguen fuente y PDF actualizado. Distingan \textbf{P1} y \textbf{P2}. Enlacen evidencias, decisiones, versiones y resultados por códigos o referencias internas, con fuentes y fechas reales.

Distingan \concept{observaciones y comentarios}, \concept{interpretaciones}, \concept{supuestos} e \concept{ideas}; identifiquen simulaciones y límites. \instruction{Usen IA solo para código LaTeX, formato y organización de contenido propio.} No inventen evidencias.

Actualicen la \textbf{bitácora} por sesión: fecha, aportes, evidencias, compromisos, responsables, plazos, estados y acuerdos. Las guías admiten proyectos físicos, digitales, comunicativos, de servicios o procedimentales. Adapten medios y extensión a la evidencia necesaria; los formatos sugeridos son opciones.

\section*{C5 de P2. Apliquen y justifiquen decisiones en el prototipo}
\begin{enumerate}[series=pasos]
\item \textbf{Identifiquen necesidades de diseño.} Recorran V1 y consulten los H y las decisiones previas. Localicen una dificultad o requisito pertinente de diseño, funcionalidad, accesibilidad o experiencia de usuario. Expliquen para quién importa y qué evidencia lo respalda; si es una hipótesis del equipo, declárenla.
\item \textbf{Decidan y construyan.} Elijan soluciones pertinentes y aplíquenlas mientras continúan la construcción. Expliquen por qué cada decisión responde al requisito y qué condiciones o alternativas consideraron. Pueden intervenir componentes, instrucciones, secuencias, interacciones o medios de acceso según el proyecto. Coordinen recursos y tareas con C6 a medida que trabajan.
\item \textbf{Hagan visible la aplicación.} Identifiquen la versión resultante como V2 (o con el código que corresponda a su historial). Localicen las decisiones en fotos, capturas, secuencias, registros de ejecución u otros soportes. Vinculen cada decisión con el principio, la necesidad o dificultad, su fundamento y el paso de R1 afectado. Señalen qué observarían en C7 para comprobar su efecto; no afirmen todavía que mejoró el uso.
\end{enumerate}

\section*{Entregables de C5}
\begin{itemize}
\item \textbf{Registro de decisiones aplicadas y justificadas en C5.} Enlacen cada principio pertinente con la dificultad o requisito, su fuente, la decisión y su fundamento, la ubicación en la versión construida y el paso de R1 afectado. Indiquen la señal observable para comprobar su efecto en C7.
\item \textbf{Prototipo vigente y soportes de aplicación.} Mantengan disponible la versión identificada y evidencias que permitan localizar las decisiones implementadas. Remitan a V1 en C4 sin volver a presentar su construcción como evidencia nueva; C5 debe mostrar la aplicación razonada durante el desarrollo.
\end{itemize}

\newpage
\section*{C6 de P2. Organicen, ejecuten y sigan la construcción}
\begin{enumerate}[resume=pasos]
\item \textbf{Coordinen el trabajo real.} Retomen PL1 y acuerden tareas de construcción vinculadas a las decisiones de C5. Precisen responsables, roles, recursos disponibles y necesarios, acceso, tiempos y dependencias. Comprueben su disponibilidad y distribuyan el trabajo según lo que requiere el avance.
\item \textbf{Ejecuten y registren el avance.} Realicen las tareas y actualicen su estado con fecha: terminada, en curso, bloqueada o pendiente, según corresponda. Indiquen qué se hizo, quién lo hizo y dónde se ve el resultado en el prototipo o sus registros. No marquen como realizada una tarea solo porque fue asignada. Conserven la planificación inicial para compararla con la ejecución.
\item \textbf{Ajusten la organización durante la construcción.} Registren obstáculos y cambios de recursos, roles, secuencia o plazos; expliquen su motivo y efecto en el avance. Relacionen los cambios con decisiones de C5 y acuerdos de bitácora. Si no fue necesario reorganizar, justifiquen cómo el seguimiento confirmó que podían continuar.
\item \textbf{Cierren el estado disponible para C7.} Identifiquen la versión que podrán probar, la función o interacción disponible y sus límites o simulaciones. Señalen qué recursos de prueba están confirmados y qué falta gestionar. Preparar la prueba puede figurar en PL1, pero no reemplaza la evidencia de construcción ejecutada.
\end{enumerate}

\section*{Entregables de C6}
\begin{itemize}
\item \textbf{PL1 con historial de seguimiento y evidencia de ejecución.} Incluyan tareas, responsables y roles, recursos, plazos, estados fechados y avance realizado, con referencias que permitan localizar los resultados en el prototipo o sus registros. Conserven la planificación inicial: \textbf{un plan sin avance documentado no evidencia por sí solo C6}.
\item \textbf{Balance de organización y enlaces a la bitácora.} Contrasten lo previsto con lo realizado; expliquen los ajustes, sus motivos y efectos, o cómo el seguimiento confirmó que podían continuar sin reorganizar. Relacionen los acuerdos con las decisiones de C5, sin duplicar V1 de C4 ni sustituir la justificación de C5 por una lista de tareas.
\end{itemize}

\section*{Resultado y continuidad}
Dejen disponibles el prototipo identificado (V2 o versión vigente), su evidencia de funcionamiento, R1 con cambios si corresponde, las decisiones de C5 y sus señales observables, PL1 actualizado y la bitácora. En C7 elegirán una prueba pertinente con usuarios, pares o docentes; en C8 aplicarán ajustes sustentados en los resultados y limitaciones.

\section*{Pregunta de reflexión}
¿Cómo influyó la manera de organizar el trabajo en las decisiones que lograron aplicar al prototipo, según los avances, obstáculos y acuerdos registrados por el equipo?

Incorporen su respuesta al documento acumulativo y susténtenla en la experiencia y las evidencias del equipo.

\end{document}
